\documentclass[11pt]{article}
\usepackage{geometry,graphicx,booktabs,xcolor,hyperref,amsmath}
\geometry{margin=1in}
\definecolor{pwmblue}{HTML}{2E5FA1}
\definecolor{pwmred}{HTML}{C93C3C}
\definecolor{pwmgreen}{HTML}{3D9E3F}

\title{\textbf{CT QC Copilot: Clinical Metrics Summary}\\[4pt]
\large Extracted from ``Physics World Models for Computational Imaging''\\
Prepared for Prof.\ Steve Jiang}
\author{Chengshuai Yang \and Xin Yuan}
\date{February 2026}

\begin{document}
\maketitle
\thispagestyle{empty}

\section*{Overview}

The CT QC Copilot translates the Triad Decomposition framework to clinical CT quality assurance. It implements ACR CT accreditation standards with automated metric computation, statistical process control for drift detection, and workflow efficiency optimization.

\section{Gate Mapping: Research $\to$ Clinical CT}

\begin{table}[htbp]
\centering
\caption{Mapping of Triad Decomposition gates to clinical CT QC failure modes.}
\begin{tabular}{lll}
\toprule
\textbf{Gate} & \textbf{Research Interpretation} & \textbf{Clinical CT Interpretation} \\
\midrule
Gate~1 & Information deficiency       & Protocol design inadequacy \\
       & (null space, compression)     & (insufficient projections, FOV truncation) \\
\midrule
Gate~2 & Carrier budget               & Dose budget \\
       & (SNR, photon count)          & (noise floor vs.\ diagnostic requirement) \\
\midrule
\textbf{Gate~3} & \textbf{Operator mismatch}    & \textbf{Scanner calibration drift} \\
       & (mask shift, PSF error)      & (HU drift, CoR offset, gain variation) \\
\bottomrule
\end{tabular}
\end{table}

\noindent\textbf{Key finding:} Gate~3 dominates in clinical practice---most QA failures trace to calibration drift, not protocol or dose issues.  This is consistent with the research findings across all 7 validated modalities.

\section{ACR Metric Validation}

Nine ACR-aligned QC metrics computed automatically from CT phantom (Gammex 464) scans on a GE Revolution Apex scanner (120~kVp, 200~mAs, 5~mm axial).

\begin{table}[htbp]
\centering
\caption{ACR metric agreement: CT QC Copilot vs.\ console gold standard.}
\begin{tabular}{lrrrr}
\toprule
\textbf{Metric} & \textbf{Copilot} & \textbf{Console} &
  \textbf{Deviation} & \textbf{\% of Tolerance} \\
\midrule
CT\# Water (HU)          & 0.3     & 0.2     & 0.1   & 2\% \\
CT\# Bone (HU)           & 952     & 953     & 1.0   & 5\% \\
CT\# Air (HU)            & $-$997  & $-$998  & 1.0   & 5\% \\
CT\# Acrylic (HU)        & 120     & 121     & 1.0   & 7\% \\
CT\# Polyethylene (HU)   & $-$96   & $-$96   & 0.0   & 0\% \\
Geometric accuracy (mm)  & 0.15    & 0.10    & 0.05  & 2.5\% \\
Slice thickness (mm)     & 5.08    & 5.02    & 0.06  & 4\% \\
Uniformity (HU)          & 1.2     & 1.1     & 0.1   & 2\% \\
Noise std.\ dev.\ (HU)  & 4.6     & 4.5     & 0.1   & 10\% \\
Spatial resolution (lp/cm) & 6.2   & 6.0     & 0.2   & 4\% \\
\bottomrule
\end{tabular}
\end{table}

\noindent\textbf{Summary:} $\leq 1.2$\,HU agreement for CT number accuracy; $\leq 0.15$\,mm for geometric accuracy.  All metrics within $\leq 10\%$ of clinical tolerance bands.

\section{Statistical Drift Detection (30-Scanner Fleet)}

Statistical process control (SPC) using five Western Electric rules, evaluated on a 30-scanner simulated fleet over 12 monthly measurement cycles.

\subsection{Drift Scenarios Tested}
\begin{itemize}
  \item Noise drift: 2 scanners (gradual noise standard deviation increase)
  \item Uniformity drift: 1 scanner (progressive uniformity degradation)
  \item CT number drift: 1 scanner (systematic HU calibration shift)
\end{itemize}

\subsection{Detection Performance}

\begin{table}[htbp]
\centering
\caption{Drift detection performance summary.}
\begin{tabular}{lr}
\toprule
\textbf{Metric} & \textbf{Value} \\
\midrule
Sensitivity              & 100\% (4/4 drifting scanners) \\
95\% CI (sensitivity)    & 39.8\%--100\% \\
Specificity              & 100\% (0/26 false positives) \\
95\% CI (specificity)    & 86.8\%--100\% \\
Early warning lead time  & 3--6 months before ACR threshold \\
False positive rate      & 0\% \\
\bottomrule
\end{tabular}
\end{table}

\section{Workflow Efficiency}

\begin{table}[htbp]
\centering
\caption{Per-scanner QC workflow time comparison: manual physicist workflow vs.\ CT QC Copilot.}
\begin{tabular}{lrr}
\toprule
\textbf{Stage} & \textbf{Manual (min)} & \textbf{Copilot (min)} \\
\midrule
DICOM ingestion          & 5.0   & 0.01 \\
Metric computation       & 25.0  & 0.01 \\
Threshold evaluation     & 5.0   & $<$0.01 \\
Trend analysis           & 10.0  & $<$0.01 \\
Report generation        & 15.0  & 0.02 \\
Physicist review/sign-off & 7.0  & 4.0 \\
\midrule
\textbf{Total}           & $\mathbf{67 \pm 12}$ & $\mathbf{4.2 \pm 0.8}$ \\
\bottomrule
\end{tabular}
\end{table}

\subsection{Annual Impact (30-Scanner Fleet)}

\begin{table}[htbp]
\centering
\caption{Annual workflow savings projection.}
\begin{tabular}{lr}
\toprule
\textbf{Metric} & \textbf{Value} \\
\midrule
Time reduction per scanner         & 94\% \\
Physicist-hours saved annually     & 377 hours \\
FTE reallocation                   & 0.18 FTE \\
Reproducibility                    & Bit-exact (SHA-256 verified) \\
Inter-analyst variability          & Eliminated \\
\bottomrule
\end{tabular}
\end{table}

\section{Artifact Signatures}

Six artifact signatures computed by the Copilot correspond directly to OperatorGraph mismatch parameter families:

\begin{table}[htbp]
\centering
\caption{CT artifact signatures and their corresponding mismatch parameters.}
\begin{tabular}{lll}
\toprule
\textbf{Artifact} & \textbf{Mismatch Parameter} & \textbf{Clinical Impact} \\
\midrule
Ring artifacts      & Center-of-rotation offset  & Diagnostic confusion \\
Cupping artifacts   & Beam hardening coefficient & CT\# inaccuracy \\
Streak artifacts    & Metal/motion               & Obscured anatomy \\
HU drift            & Detector gain calibration  & Quantitative error \\
Noise ratio change  & Tube current / filtration  & Dose optimization \\
Geometric distortion & Gantry alignment          & Treatment planning error \\
\bottomrule
\end{tabular}
\end{table}

\section{CT Correction Results}

From the main paper validation:
\begin{itemize}
  \item \textbf{Mismatch type}: Sub-degree center-of-rotation offset
  \item \textbf{Scenario II PSNR}: 13.41\,dB (severe ring artifacts)
  \item \textbf{Scenario III PSNR}: 24.09\,dB (corrected)
  \item \textbf{Correction gain}: $+10.68$\,dB
  \item \textbf{Solver}: FBP with Ram-Lak filter
\end{itemize}

\section*{Opportunities for Collaboration}

\begin{enumerate}
  \item \textbf{Prospective clinical validation}: Deploy CT QC Copilot on real clinical scanner fleet at UT Southwestern with physical ACR phantoms, replacing simulated fleet validation.
  \item \textbf{MRI clinical validation}: Validate coil sensitivity mismatch correction on clinical MRI data (fastMRI + institutional data).
  \item \textbf{Regulatory pathway}: Guide FDA 510(k) or De Novo pathway for CT QC Copilot as clinical decision support software.
  \item \textbf{Multi-institution study}: Extend drift detection validation across institutions with diverse scanner manufacturers (GE, Siemens, Philips, Canon).
\end{enumerate}

\end{document}
